% Generated draft: requires mathematical and source review.
\documentclass[11pt]{article}
\usepackage[margin=25mm]{geometry}
\usepackage{fontspec}
\setmainfont{Latin Modern Roman}
\usepackage{amsmath,amssymb,mathtools,unicode-math}
\setmathfont{Latin Modern Math}
\usepackage{enumitem,needspace,xurl,hyperref}
\setlength{\parindent}{0pt}
\setlength{\parskip}{5pt}
\setlength{\emergencystretch}{3em}
\newcommand{\N}{\mathbb{N}}
\newcommand{\Z}{\mathbb{Z}}
\newcommand{\Q}{\mathbb{Q}}
\newcommand{\R}{\mathbb{R}}
\newcommand{\C}{\mathbb{C}}
\newcommand{\F}{\mathbb{F}}
\newcommand{\A}{\mathbb{A}}
\newcommand{\PP}{\mathbb{P}}
\newcommand{\E}{\mathbb{E}}
\newcommand{\eps}{\varepsilon}
\newcommand{\dd}{\,\mathrm{d}}
\begin{document}
\section*{UnsolvedMath ID: 30002120}
\textbf{Source problem code:} catalogue entry
\subsection{Definitions and mathematical statement}
Let \(m\geq 1\) be an integer and let \(0<s<1\). Consider a lower-semibounded fractional Schrödinger operator
\[
H=(-\Delta)^s+V\quad\text{on }L^2(\R^m),
\]
where \(V\) is real. Here \((-\Delta)^s\) denotes the fractional Laplacian on the whole space. The supplied record does not specify hypotheses on \(V\), an operator domain, or a self-adjoint realization; the spectral terminology below requires a suitable realization to be fixed.

Suppose that the \(k\)th eigenvalue below the essential spectrum exists, and denote it by \(\lambda_k\). Let \(\psi_k\) be a nonzero, continuous, real-valued eigenfunction satisfying
\[
\psi_k\in\mathcal D(H),\qquad H\psi_k=\lambda_k\psi_k.
\]
The source does not specify how eigenvalue multiplicities are counted when assigning the index \(k\).

Define the nodal domains of \(\psi_k\) to be the connected components, in the Euclidean topology of \(\R^m\), of the open set
\[
\Omega_{\psi_k}=\{x\in\R^m:\psi_k(x)\neq 0\}.
\]
Write
\[
\nu(\psi_k)=\#\{D:D\text{ is a connected component of }\Omega_{\psi_k}\},
\]
allowing \(\nu(\psi_k)=\infty\).

\textbf{Problem.} Bound \(\nu(\psi_k)\) in terms of the eigenvalue index \(k\), for eigenfunctions satisfying the preceding assumptions. The supplied formulation does not state whether the bound must be uniform in \(m\), \(s\), \(V\), and the operator realization, or whether dependence on some of these data is permitted. Consequently, it does not determine a single fully quantified bound to be established.

\textbf{Formulation qualification.} The registry's original statement concerns eigenfunctions of a Dirichlet Laplacian or a comparable elliptic operator. Its corrected statement instead concerns the fractional Schrödinger operator on \(\R^m\) above. These are different settings. This draft follows the corrected statement without importing a Dirichlet boundary condition. The registry labels the correction \texttt{corrected\_verified} and attributes it to the official report; that source-verification claim is supplied metadata, not an independent verification performed for this draft.
\subsection{Short English statement}
For a continuous real eigenfunction of a fractional Schrödinger operator on \(\R^m\), count the connected regions on which the eigenfunction is nonzero and seek an upper bound in terms of its eigenvalue index. The potential assumptions, operator realization, multiplicity convention, and permitted parameter dependence of the bound require clarification. This is a statement-only draft requiring human mathematical review.
\subsection{Sources}
\begin{itemize}
\item Local UnsolvedMath registry, supplied complete record: permanent ID \(30002120\), source problem code \texttt{catalogue entry}, title \emph{Nodal Counts for Fractional Schrödinger Eigenfunctions}; collection: \emph{Oberwolfach Reports Open Problems}. No individual proposer is supplied. The proposed year is recorded as 2012.
\item Supplied report citation: \emph{Geometric Aspects of Spectral Theory (2013)}, Oberwolfach Reports, pp.\ 2013--2076, DOI \texttt{10.4171/owr/2012/33}. \url{https://doi.org/10.4171/owr/2012/33}. The supplied background describes Oberwolfach Reports as documenting mathematical workshops and their open-problem discussions.
\item Supplied literature source, labeled ``Hur Johnson Martin'': \url{https://doi.org/10.19086/da.2102}. No further bibliographic details are supplied.
\item Supplied literature source, labeled ``OWR report'': \url{https://ems.press/content/serial-article-files/46403}.
\item Supplied literature assessment dated \textbf{2026-08-21}: the registry labels the problem open and reports that no general index-versus-nodal-domain theorem is known for fractional Schrödinger eigenfunctions, while one-dimensional and other partial oscillation bounds are available. This is supplied dated curation, described in the record as literature triage derived from OWR curation; it has not been independently verified for this draft. The record expressly requests independent verification of time-sensitive status and cited claims.
\end{itemize}
\end{document}
